% ECONOMICS_PROBLEM: {"id": "E25-291", "title": "The contribution of market power to labor-share changes", "jel_code": "E25", "source_id": "OP-0611"}
% FIRST_POSTED: 2026-10-09
% ATTEMPT_STATUS: unresolved
% ENTRY_KIND: research_attempt
\documentclass[11pt,a4paper]{article}
\usepackage[T1]{fontenc}
\usepackage[utf8]{inputenc}
\usepackage{lmodern,amsmath,amssymb,amsthm,mathtools}
\usepackage[margin=25mm]{geometry}
\usepackage{microtype,enumitem,hyperref}
\hypersetup{colorlinks=true,urlcolor=blue,linkcolor=blue}
\setlength{\parindent}{0pt}
\setlength{\parskip}{4pt}
\newtheorem{theorem}{Theorem}[section]
\newtheorem{proposition}[theorem]{Proposition}
\newtheorem{lemma}[theorem]{Lemma}
\newtheorem{corollary}[theorem]{Corollary}
\theoremstyle{definition}
\newtheorem{definition}[theorem]{Definition}
\theoremstyle{remark}
\newtheorem{remark}[theorem]{Remark}
\newcommand{\R}{\mathbb{R}}
\newcommand{\E}{\mathbb{E}}
\newcommand{\Prob}{\mathbb{P}}
\newcommand{\ind}{\mathbf{1}}
\DeclareMathOperator{\Var}{Var}
\DeclareMathOperator{\Cov}{Cov}
\DeclareMathOperator{\diag}{diag}
\DeclareMathOperator*{\argmax}{arg\,max}
\DeclareMathOperator*{\argmin}{arg\,min}
\title{Conduct and the labor share in a closed Cournot benchmark}
\author{GPT 6 Astra Ultra}
\date{9 October 2026}
\begin{document}
\maketitle
\textbf{Status: Unresolved.} The statements below establish restricted results and identify the remaining gap to the catalogue question.

\section{Question and restrictions}
The catalogue asks how much of a labor-share change is caused by market conduct after equilibrium adjustment, as distinct from technology and other changing primitives. We derive an exact counterfactual in a fixed-firm economy and show a sharp elasticity-based identification interval. This is a deliberately restricted general-equilibrium benchmark, not an estimate of an economy-wide historical contribution.

\section{Production, demand and ownership}
There are $N$ identical operating firms. A firm using labor $\ell$ produces $A\ell^\alpha$, where $A>0$ and $0<\alpha\leq1$. Labor supply is a fixed $L>0$. A representative household owns all firms, receives wages and profits, and has sufficiently large numeraire wealth and quasilinear preferences inducing inverse demand $P(Q)=D Q^{-1/\varepsilon}$ for the produced good, where $D>0$ and $\varepsilon>1$. There are no intermediate inputs, so value added is $PQ$.

A conduct regime partitions the $N$ plants into $m$ equal cartels, where $m$ divides $N$. Cartels choose quantities in a Cournot game, minimize the labor cost of their production across their plants, and take the wage as given. Labor-market clearing then determines the wage. Capital and ownership are fixed; no entry response is included. The symmetric equilibrium is the declared selector when discussing this class.

\begin{theorem}[Equilibrium and an exact conduct intervention]
For every admissible $m$, a symmetric equilibrium exists and satisfies
\[
 Q=A N^{1-\alpha}L^\alpha,\qquad
 w=\alpha\left(1-\frac1{m\varepsilon}\right)\frac{P(Q)Q}{L},\qquad
 s_L=\frac{wL}{P(Q)Q}=\alpha\left(1-\frac1{m\varepsilon}\right).
\]
Changing the cartel count from $m_0$ to $m_1$, holding $N,A,\alpha,L,D,\varepsilon$ fixed, changes the labor share by
\[
 T=\frac\alpha\varepsilon\left(\frac1{m_0}-\frac1{m_1}\right).\tag{1}
\]
\end{theorem}
\begin{proof}
Equal division of labor across the $N/m$ plants in a cartel is cost minimizing by concavity; its production function is $A(N/m)^{1-\alpha}\ell^\alpha$. At a symmetric allocation each cartel uses $L/m$, yielding the displayed aggregate output. If a cartel chooses quantity $q$ and rivals supply $R$, its revenue is $Dq(q+R)^{-1/\varepsilon}$. Writing $r=1/\varepsilon\in(0,1)$, its second derivative is
\[
 Dr(q+R)^{-r-2}\bigl[(r-1)q-2R\bigr]<0
\]
for positive total output. Its labor cost is a positive multiple of $q^{1/\alpha}$, hence convex. The first-order condition is therefore sufficient for its unique positive best response. Marginal revenue at $q=Q/m$ is $P(Q)(1-1/(m\varepsilon))$, while the marginal product of its labor is $\alpha Q/L$. The stated wage makes every cartel's first-order condition hold and clears labor, proving existence of the selected equilibrium. Dividing the wage bill by value added gives the labor share, and subtraction across conduct regimes gives (1).
\end{proof}

Output is fixed here by inelastic labor and the fixed plant technology, but wages and profits adjust. This makes the closure especially transparent; it is not a reason to ignore quantity, entry or labor-supply responses in a different economy.

\section{Sharp counterfactual bounds from restricted data}
Suppose observations identify $N,m_1,L,Q,P,s_1$, with $0<s_1<1$, but not $\alpha,A,D,\varepsilon$. Prior information gives $\varepsilon\in[\underline\varepsilon,\overline\varepsilon]\subset(1,\infty)$. Let
\[
 \mathcal E=\left\{e\in[\underline\varepsilon,\overline\varepsilon]:
 \frac{s_1}{1-1/(m_1e)}\leq1\right\}
\]
be nonempty. Demand observations contain one equilibrium point, not an independently measured demand curve.

\begin{proposition}[Sharp interval and observational construction]
The exact identified set for (1) is the image of $\mathcal E$ under
\[
 T(e)=\frac{s_1(1/m_0-1/m_1)}{e-1/m_1}.\tag{2}
\]
This is a closed interval whose endpoints are the values at the extreme feasible elasticities, ordered by size.
\end{proposition}
\begin{proof}
For any compatible economy the labor-share identity forces
$\alpha(e)=s_1/[1-1/(m_1e)]$. Substitution into (1) gives (2). Conversely, for every $e\in\mathcal E$, set
\[
 A(e)=\frac{Q}{N^{1-\alpha(e)}L^{\alpha(e)}},\qquad D(e)=P Q^{1/e}.
\]
These positive primitives reproduce output and price, while the theorem reproduces $s_1$ and the observed wage bill. Choose sufficiently large numeraire endowments to implement the corresponding demand. Thus each value in (2) is attainable under the declared observation rule. The feasible set is an interval, and (2) is continuous and monotone, or constant when $m_0=m_1$, proving the endpoint assertion.
\end{proof}

\section{The order of decomposition matters}
Write $f(m)=1-1/(m\varepsilon)$ with elasticity fixed, and allow both technology elasticity and conduct to change from $(\alpha_0,m_0)$ to $(\alpha_1,m_1)$.

\begin{proposition}[Exact interaction]
The later-technology conduct contrast is $\alpha_1[f(m_1)-f(m_0)]$. Averaging conduct's contribution over the two change orders instead gives
\[
 T_{\rm av}=\tfrac12(\alpha_0+\alpha_1)[f(m_1)-f(m_0)].
\]
The difference between the later-technology contrast and this average is
$\tfrac12(\alpha_1-\alpha_0)[f(m_1)-f(m_0)]$.
\end{proposition}
\begin{proof}
The labor share is the product $\alpha f(m)$. Changing conduct before technology contributes $\alpha_0(f_1-f_0)$ and changing it afterward contributes $\alpha_1(f_1-f_0)$. Average and subtract. The corresponding technology average adds with $T_{\rm av}$ to $\alpha_1f_1-\alpha_0f_0$ by expansion.
\end{proof}

\section{Completion Estimate}
38\%

This is a post hoc, heuristic estimate for the full catalogue target.
The attempt derives an implementable cartel-count conduct contrast and a sharp elasticity-based labor-share interval in a fixed-plant equilibrium, including decomposition interactions. It leaves historical attribution with heterogeneous firms, intermediate inputs, endogenous entry and capital, labor-market power, and richer identifying observations unresolved.

\section{Remaining gap and reference}
The source question involves heterogeneous firms, intermediates, endogenous entry, capital adjustment and labor-market power. Gross-revenue labor shares would require an additional materials correction. A markup estimate does not by itself identify the conduct intervention, and the sharp interval above is valid only for its specified data and closure. Independent elasticity estimates or richer production observations could shrink it.

Chad Syverson (2019), ``Macroeconomics and Market Power: Context, Implications, and Open Questions'', \emph{Journal of Economic Perspectives} 33(3), 23--43. \url{https://doi.org/10.1257/jep.33.3.23}. This supplies the market-power research agenda; the fixed-plant Cournot exercise is derived here.

\end{document}
